For message decoding, assume $t\mid q$ and $\Delta=q/t$. Work with fixed integer coefficient lifts of the public multiplier $\Lambda$ when bounding noise. Equality of ciphertext phases below does not assert equality of fresh-encryption distributions.

\textbf{- Reference:} 
\href{https://www.zama.ai/post/tfhe-deep-dive-part-2}{TFHE Deep Dive - Part II - Encodings and linear leveled operations}~\cite{tfhe-2}

$ $

Suppose we have a GLWE ciphertext \textsf{ct}:

$\textsf{ct} = \textsf{GLWE}_{S, \sigma}(\Delta M + E) = ( A_0, A_1, \ldots, A_{k-1}, B) \in \mathcal{R}_{\langle n,q \rangle}^{k + 1}$

$ $

\noindent and a new plaintext polynomial $\Lambda$ as follows: 

$\Lambda = \sum\limits_{i=0}^{n-1}(\Lambda_i \cdot X^i) \in \mathcal{R}_{\langle n, q \rangle}$

$ $

\noindent Let's define the following ciphertext-to-plaintext multiplication operation:

$\Lambda \cdot \textsf{ct} = (\Lambda \cdot A_0, \Lambda \cdot A_1, \ldots, \Lambda \cdot A_{k-1}, \Lambda \cdot B)$

$ $

\noindent NTT multiplication uses $O(n\log n)$ field operations when the transform's root and invertibility conditions hold (\autoref{sec:ntt}); for instance, prime $q$ with $2n\mid q-1$ supports the negacyclic transform. An arbitrary modulus does not automatically provide such an NTT. The following algebraic identity holds independently of the multiplication algorithm:

\begin{tcolorbox}[title={\textbf{\tboxlabel{\ref*{sec:glwe-mult-plain}} GLWE Ciphertext-to-Plaintext Multiplication}}]
$\Lambda \cdot \textsf{GLWE}_{S, \sigma}(\Delta M + E)$

$= \Lambda \cdot (\{A_i^{\langle 1 \rangle}\}_{i=0}^{k-1}, \text{ } B^{\langle 1 \rangle})$

$= (\{\Lambda\cdot A_i^{\langle 1 \rangle}\}_{i=0}^{k-1}, \text{ } \Lambda \cdot B^{\langle 1 \rangle})$

$= \textsf{GLWE}_{S, \sigma}(\Delta (M \cdot \Lambda) + \Lambda\cdot E )$
\end{tcolorbox}

The resulting phase is $\Delta M\Lambda+\Lambda E$. It decodes to $M\Lambda\bmod t$ only when $\|\Lambda E\|_\infty<\Delta/2$.

$ $

%\noindent \textbf{\underline{Proof}}
\begin{myproof}
\begin{enumerate}
\item Define the following notations: \\
$A_0' = \Lambda \cdot A_0$ \\
$A_1' = \Lambda \cdot A_1$ \\
$\vdots$ \\
$A_{k-1}' = \Lambda \cdot A_{k-1}$ \\
$E' = \Lambda \cdot E$ \\
$B' = \Lambda \cdot B$ \\
\item Derive the following: \\
$B' = \Lambda \cdot B$ \\
$= \Lambda \cdot (\sum\limits_{i=0}^{k-1}{(A_i \cdot S_i)} + \Delta \cdot M + E)$ 
$= \sum\limits_{i=0}^{k-1}{(\Lambda \cdot A_i \cdot S_i)} + \Delta \cdot \Lambda \cdot M + \Lambda \cdot E$   \\ \textcolor{red}{ $\rhd$ by the distributive property of a polynomial ring} \\
$= \sum\limits_{i=0}^{k-1}{((\Lambda \cdot A_i) \cdot S_i)} + \Delta \cdot (\Lambda \cdot M) + (\Lambda \cdot E)$  \\
$= \sum\limits_{i=0}^{k-1}{(A_i' \cdot S_i)} + \Delta \cdot (\Lambda \cdot M) + (E')$ \\
\item Since $B' = \sum\limits_{i=0}^{k-1}{(A_i' \cdot S_i)} + \Delta \cdot (\Lambda \cdot M) + (E')$, 

$(A_0', A_1', \ldots, A_{k-1}'
, B')$ form the ciphertext $\textsf{GLWE}_{S, \sigma}(\Delta \cdot \Lambda \cdot M+E^{\prime})$.
\item Thus, \\
$\Lambda \cdot \textsf{GLWE}_{S, \sigma}(\Delta M + E)$ \\
$ = (\Lambda \cdot A_0, \text { } \Lambda \cdot A_1, \ldots, \Lambda \cdot A_{k-1}, \text { } \Lambda \cdot B)$ \\
$ = ( \{A'_{i}\}_{i=0}^{k-1}, \text { } \Lambda \cdot B)$ \\
$= \textsf{GLWE}_{S, \sigma}(\Delta (M \cdot \Lambda) + \Lambda \cdot E)$

%\begin{flushright}
%\qedsymbol{} 
%\end{flushright}
\end{enumerate}
\end{myproof}

Secret-key subtraction removes the mask products exactly, but rounding removes $E'=\Lambda E$ only within the decoding bound. In the negacyclic ring, each output coefficient is a signed sum of at most $n$ products, so
\[
\|\Lambda E\|_\infty\le\|\Lambda\|_1\|E\|_\infty\le n\|\Lambda\|_\infty\|E\|_\infty.
\]
Here $\|\Lambda\|_1$ is the sum of the absolute values of its lifted coefficients.

Repeated multiplications can increase the noise bound until decoding fails. Decryption reads the result but does not refresh an existing ciphertext. For large multipliers, a precomputed GLev representation can give a better bound through gadget decomposition (\autoref{subsec:gadget-decomposition}), as below.

\subsection{Gadget Decomposition for Noise Suppression}
\label{subsubsec:gadget-decomposition-noise-suppression}

Here use the exact descending gadget $g_i=q/\beta^i$ with $q=\beta^l$ and unsigned digits in $\{0,\ldots,\beta-1\}$. The precomputed input is $\textsf{GLev}^{\beta,l}_{S,\sigma}(\Delta M)$, not an arbitrary single GLWE ciphertext. Such a low-noise GLev representation is additional encrypted input and cannot be assumed to appear for free from a ciphertext after previous operations.

In the ciphertext-to-plaintext multiplication $\Lambda \cdot \textsf{GLWE}_{S, \sigma}(\Delta M)$, the noise $E$ grows to $E' = \Lambda \cdot E$. To limit this noise growth, we introduce a technique based on decomposing $\Lambda$ (\autoref{subsec:number-decomp}) and a GLev encryption (\autoref{subsec:glev-enc}) of $M$ as follows:


$\Lambda = \Lambda_1 \dfrac{q}{\beta^1} + \Lambda_2 \dfrac{q}{\beta^2} + \cdots + \Lambda_l \dfrac{q}{\beta^l} \longrightarrow \textsf{Decomp}^{\beta, l}(\Lambda) = (\Lambda_1, \Lambda_2, \cdots, \Lambda_l)$

$ $

$\textsf{GLev}_{S, \sigma}^{\beta, l}(\Delta M) = \Bigg\{ \textsf{GLWE}_{S, \sigma}\left(\Delta M \dfrac{q}{\beta^1} + E_1\right), \textsf{GLWE}_{S, \sigma}\left(\Delta M \dfrac{q}{\beta^2} + E_2\right), \cdots \textsf{GLWE}_{S, \sigma}\left(\Delta M \dfrac{q}{\beta^l} + E_l\right) \Bigg\}$

$ $


We will encrypt the plaintext $M$ as $\textsf{GLev}_{S, \sigma}^{\beta, l}(\Delta M)$ instead of $\textsf{GLWE}_{S, \sigma}(\Delta M)$, and compute $\textsf{Decomp}^{\beta, l}(\Lambda) \cdot \textsf{GLev}_{S, \sigma}^{\beta, l}(\Delta M)$ instead of $\Lambda \cdot \textsf{GLWE}_{S, \sigma}(\Delta M)$. Notice that the results of both computations are the same as follows:

$\textsf{Decomp}^{\beta, l}(\Lambda) \cdot \textsf{GLev}_{S, \sigma}^{\beta, l}(\Delta M)$

$= (\Lambda_1, \Lambda_2, \cdots, \Lambda_l) \cdot \left (\textsf{GLWE}_{S, \sigma}\left(\dfrac{q}{\beta} \Delta M + E_1\right), \text{ } \textsf{GLWE}_{S, \sigma}\left(\dfrac{q}{\beta^2} \Delta M + E_2\right), \text{ } \cdots, \text{ } \textsf{GLWE}_{S, \sigma}\left(\dfrac{q}{\beta^l} \Delta M + E_l\right) \right )$

$= \Lambda_1\cdot\textsf{GLWE}_{S, \sigma}\left(\dfrac{q}{\beta} \Delta M + E_1\right) +  \Lambda_2\cdot\textsf{GLWE}_{S, \sigma}\left(\dfrac{q}{\beta^2} \Delta M + E_2\right) + \cdots + \Lambda_l\cdot\textsf{GLWE}_{S, \sigma}\left(\dfrac{q}{\beta^l} \Delta M + E_l\right)$

$= \textsf{GLWE}_{S, \sigma}\left(\Lambda_1\cdot\dfrac{q}{\beta}\Delta M + \Lambda_1 E_1\right) +\textsf{GLWE}_{S, \sigma}\left(\Lambda_2\cdot\dfrac{q}{\beta^2}\Delta M + \Lambda_2 E_2\right)+ \cdots + \textsf{GLWE}_{S, \sigma}\left(\Lambda_l\cdot\dfrac{q}{\beta^l}\Delta M + \Lambda_l E_l\right)$

$= \textsf{GLWE}_{S, \sigma}\left(\Lambda_1\cdot\dfrac{q}{\beta} \Delta M + \Lambda_2\cdot\dfrac{q}{\beta^2} \Delta M + \cdots + \Lambda_l\cdot\dfrac{q}{\beta^l} \Delta M+E_{\textit{all}}\right)$

$= \textsf{GLWE}_{S, \sigma}\left(\left(\Lambda_1\cdot\dfrac{q}{\beta} + \Lambda_2\cdot\dfrac{q}{\beta^2} + \cdots + \Lambda_l\cdot\dfrac{q}{\beta^l}\right)\cdot \Delta M  + E_{\textit{all}} \right)$ \textcolor{red}{ $\rhd$ where $E_{\textit{all}} = \sum\limits_{i=1}^l \Lambda_iE_i$}

$= \textsf{GLWE}_{S, \sigma}\left(\Lambda \cdot \Delta M + E_{\textit{all}}\right)$ \textcolor{red}{ $\rhd$ whose decryption is $\Lambda\cdot M$}

$ $

For exact decomposition, the final error is $E_{\textit{all}}=\sum_{i=1}^l\Lambda_iE_i$. If $\|E_i\|_\infty\le B_{\mathrm{g}}$, the unsigned digits give
\[
\|E_{\textit{all}}\|_\infty\le\sum_{i=1}^l\|\Lambda_i\|_1\|E_i\|_\infty\le nl(\beta-1)B_{\mathrm{g}}.
\]
A $\beta/2$ digit bound would require an explicitly balanced decomposition instead. This can improve the worst-case bound compared with a large multiplier, but does not guarantee a smaller realized error for every input. Both computations decode to the same message only if their respective error bounds hold. For truncated decomposition $\Lambda\equiv\sum_i\Lambda_i g_i+\rho\pmod q$, the final phase instead has error $\sum_i\Lambda_iE_i-\rho\Delta M$; this residual must also be bounded. Figure~\ref{fig:decomp2} illustrates the intended bound reduction.

\begin{figure}[h!]
    \centering
  \includegraphics[width=0.8\linewidth]{figures/decomp2.pdf}
  \caption{Noise reduction in ciphertext-to-plaintext multiplication by gadget decomposition.}
  \label{fig:decomp2}
\end{figure}

\subsubsection{Discussion}
\label{subsubsec:glwe-mult-plain-discuss}


Nevertheless, the decomposition technique is still very useful: for GLWE key-switching (\autoref{sec:glwe-key-switching}), we will show how to key-switch by combining decomposed mask polynomials $\textsf{Decomp}^{\beta,l}(A_i)$
with a precomputed key-switching key
$\textsf{KSK}_i=\textsf{GLev}^{\beta,l}_{S',\sigma}(S_i)$,
so gadget decomposition can be repeatedly leveraged across key-switching calls even though each individual application outputs a standard GLWE ciphertext.

Meanwhile, for the technique to repeatedly re-initialize the noise $E$ of regular ciphertexts, we will describe TFHE's noise bootstrapping technique in \autoref{subsec:tfhe-noise-bootstrapping}.
